\documentclass[aspectratio=169,11pt]{beamer}
% Shared Beamer style for practice contest solution walkthroughs.
\usepackage[utf8]{inputenc}
\usepackage[T1]{fontenc}
\usepackage{lmodern}
\usepackage{amsmath,amssymb}
\usepackage{xcolor}
\usepackage{hyperref}
\usepackage{listings}

\definecolor{CPNavy}{HTML}{172033}
\definecolor{CPBlue}{HTML}{2563EB}
\definecolor{CPGreen}{HTML}{16A34A}
\definecolor{CPOrange}{HTML}{F97316}
\definecolor{CPGray}{HTML}{64748B}
\definecolor{CPLight}{HTML}{F8FAFC}
\definecolor{CPCodeBg}{HTML}{F1F5F9}
\definecolor{CPCodeRule}{HTML}{CBD5E1}

\usetheme{default}
\usefonttheme{professionalfonts}
\setbeamertemplate{navigation symbols}{}
\setbeamersize{text margin left=0.65cm,text margin right=0.65cm}
\setbeamertemplate{itemize item}{\color{CPBlue}$\blacktriangleright$}
\setbeamertemplate{itemize subitem}{\color{CPGray}$\triangleright$}

\setbeamercolor{normal text}{fg=CPNavy,bg=white}
\setbeamercolor{frametitle}{fg=white,bg=CPNavy}
\setbeamercolor{title}{fg=white,bg=CPNavy}
\setbeamercolor{subtitle}{fg=CPLight,bg=CPNavy}
\hypersetup{colorlinks=true,urlcolor=CPBlue}

\setbeamertemplate{frametitle}{%
  \nointerlineskip%
  \begin{beamercolorbox}[wd=\paperwidth,ht=0.88cm,dp=0.22cm,leftskip=0.55cm,rightskip=0.35cm]{frametitle}
    \usebeamerfont{frametitle}\insertframetitle\hfill\small\insertframenumber/\inserttotalframenumber
  \end{beamercolorbox}%
}

\setbeamertemplate{footline}{%
  \leavevmode%
  \hbox{%
    \begin{beamercolorbox}[wd=.58\paperwidth,ht=2.4ex,dp=1ex,leftskip=.5cm]{author in head/foot}%
      \color{CPGray}\insertshorttitle
    \end{beamercolorbox}%
    \begin{beamercolorbox}[wd=.42\paperwidth,ht=2.4ex,dp=1ex,rightskip=.5cm plus1fil]{date in head/foot}%
      \hfill\color{CPGray}\insertshortauthor
    \end{beamercolorbox}%
  }%
}

\setbeamertemplate{title page}{%
  \vspace*{1.25cm}
  \begin{beamercolorbox}[wd=\paperwidth,sep=0.65cm,leftskip=0.15cm,rightskip=0.15cm]{title}
    {\color{white}\Huge\bfseries\inserttitle\par}
    \vspace{0.35cm}
    {\color{CPLight}\Large\insertsubtitle\par}
    \vspace{0.6cm}
    {\color{CPLight}\large\insertauthor\par}
  \end{beamercolorbox}
  \vfill
}

\newcommand{\sectiontag}[1]{\textbf{\color{CPBlue}#1}}
\newenvironment{tightitemize}{\begin{itemize}\small\setlength{\itemsep}{0.18cm}}{\end{itemize}}
\lstdefinestyle{cpcode}{
  language=Python,
  basicstyle=\ttfamily\tiny,
  keywordstyle=\color{CPBlue}\bfseries,
  commentstyle=\color{CPGray}\itshape,
  stringstyle=\color{CPGreen},
  backgroundcolor=\color{CPCodeBg},
  frame=single,
  framerule=0.25pt,
  rulecolor=\color{CPCodeRule},
  showstringspaces=false,
  columns=fullflexible,
  keepspaces=true,
  breaklines=true,
  breakatwhitespace=false,
  tabsize=4,
  xleftmargin=0.08cm,
  xrightmargin=0.08cm,
  aboveskip=0.08cm,
  belowskip=0.08cm
}


\title[spiderman]{Spiderman's Workout}
\subtitle{Day 3 solution walkthrough}
\author[Solution Deck]{Competitive Programming Bootcamp}
\date{}

\begin{document}

\begin{frame}[plain]
  \titlepage
\end{frame}

% BEGIN GENERATED STATEMENT INTERPRETATION
\begin{frame}{Interpret the Word Problem}
\small
\sectiontag{Statement excerpts:}
\begin{tightitemize}
  \item \textit{``starts and finishes at street level''}
  \item \textit{``minimizes the required building height''}
\end{tightitemize}

\vspace{0.18cm}
\sectiontag{Programming interpretation:}
\begin{tightitemize}
  \item Each distance must be assigned either U or D in the original order.
  \item The path may never go below height zero and must end at zero.
  \item Among valid paths, minimize the highest height ever reached.
\end{tightitemize}
\end{frame}
% END GENERATED STATEMENT INTERPRETATION







\begin{frame}{Problem Model}
\small
\sectiontag{Textbook connection:} CP4 3.5 f. DP Level 1

\vspace{0.25cm}
\sectiontag{Core technique:} DP over step index and current height

\vspace{0.25cm}
\begin{tightitemize}
  \item For each distance, Spiderman chooses U or D.
  \item Height can never become negative.
  \item He must finish at height 0 while minimizing the highest height reached.
\end{tightitemize}
\end{frame}

\begin{frame}{How to Recognize the Approach}
\begin{tightitemize}
  \item A greedy choice can trap the path later.
  \item State must include how many moves are done and the current height.
  \item Path reconstruction is required to print U/D moves.
\end{tightitemize}
\end{frame}

\begin{frame}{Algorithm}
\begin{tightitemize}
  \item Let max\_height be the sum of all distances.
  \item dp[i][h] stores the smallest maximum height after i moves ending at h.
  \item From each reachable state, try moving up and down.
  \item Update parent[i+1][new\_h] with previous height and move letter.
  \item If dp[m][0] is reachable, follow parent pointers backward to build the path.
\end{tightitemize}
\end{frame}

\begin{frame}{Walkthrough of the Key State}
\begin{tightitemize}
  \item Moving up may increase the maximum height reached.
  \item Moving down preserves the previous maximum height.
  \item The parent table turns an optimal value DP into an output-producing DP.
\end{tightitemize}
\end{frame}

\begin{frame}{Why the Algorithm Is Correct}
\begin{tightitemize}
  \item The DP considers both legal choices after every reachable prefix.
  \item For each state, it keeps the path with the smallest maximum height, which dominates worse paths to the same state.
  \item The final state (all moves, height 0) is exactly the problem's ending condition.
\end{tightitemize}
\end{frame}

\begin{frame}{Complexity and Implementation Notes}
\small
\sectiontag{Complexity:} O(MS) time and O(MS) memory, where S is the sum of distances.

\vspace{0.3cm}
\begin{tightitemize}
  \item Return IMPOSSIBLE when height 0 is unreachable after all moves.
  \item The reconstructed path is built backward and then reversed.
\end{tightitemize}
\end{frame}



% BEGIN GENERATED CODE WALKTHROUGH
\begin{frame}[fragile]{Code: Set Up DP State}
\scriptsize
\sectiontag{Source lines:} \texttt{spiderman.py:44--66}

\vspace{0.08cm}
\begin{columns}[T,totalwidth=\textwidth]
\begin{column}{0.58\textwidth}
\begin{lstlisting}[style=cpcode]
import sys

def solve(distances):
    m = len(distances)

    max_height = sum(distances)

    INF = float("inf")

    dp = [[INF] * (max_height + 1) for _ in range(m + 1)]

    parent = [[None] * (max_height + 1) for _ in range(m + 1)]

    dp[0][0] = 0
\end{lstlisting}
\end{column}
\begin{column}{0.38\textwidth}
\sectiontag{What this chunk does}
\begin{tightitemize}
  \item The maximum possible height is the sum of all distances.
  \item dp[i][h] is the lowest peak height for paths ending at h after i moves.
  \item parent stores the previous height and move so the final string can be rebuilt.
\end{tightitemize}
\end{column}
\end{columns}
\end{frame}

\begin{frame}[fragile]{Code: Try Moving Up}
\scriptsize
\sectiontag{Source lines:} \texttt{spiderman.py:68--97}

\vspace{0.08cm}
\begin{columns}[T,totalwidth=\textwidth]
\begin{column}{0.58\textwidth}
\begin{lstlisting}[style=cpcode]
for i in range(m):
    distance = distances[i]

    for height in range(max_height + 1):

        if dp[i][height] == INF:
            continue

        new_height = height + distance

        if new_height <= max_height:

            highest = max(dp[i][height], new_height)

            if highest < dp[i + 1][new_height]:
                dp[i + 1][new_height] = highest

                parent[i + 1][new_height] = (height, "U")
\end{lstlisting}
\end{column}
\begin{column}{0.38\textwidth}
\sectiontag{What this chunk does}
\begin{tightitemize}
  \item Only reachable prior heights are expanded.
  \item Moving up changes the current height and may increase the path's maximum height.
  \item The transition is kept only if it improves the best known peak for the new state.
\end{tightitemize}
\end{column}
\end{columns}
\end{frame}

\begin{frame}[fragile]{Code: Try Moving Down}
\scriptsize
\sectiontag{Source lines:} \texttt{spiderman.py:98--113}

\vspace{0.08cm}
\begin{columns}[T,totalwidth=\textwidth]
\begin{column}{0.58\textwidth}
\begin{lstlisting}[style=cpcode]
new_height = height - distance

if new_height >= 0:

    highest = dp[i][height]

    if highest < dp[i + 1][new_height]:
        dp[i + 1][new_height] = highest

        parent[i + 1][new_height] = (height, "D")
\end{lstlisting}
\end{column}
\begin{column}{0.38\textwidth}
\sectiontag{What this chunk does}
\begin{tightitemize}
  \item A downward move is legal only if the new height stays nonnegative.
  \item Moving down cannot increase the maximum height reached so far.
  \item The parent pointer records that this state came from a D move.
\end{tightitemize}
\end{column}
\end{columns}
\end{frame}

\begin{frame}[fragile]{Code: Finish and Reconstruct}
\scriptsize
\sectiontag{Source lines:} \texttt{spiderman.py:115--139}

\vspace{0.08cm}
\begin{columns}[T,totalwidth=\textwidth]
\begin{column}{0.58\textwidth}
\begin{lstlisting}[style=cpcode]
if dp[m][0] == INF:
    return "IMPOSSIBLE"

path = []
height = 0

for i in range(m, 0, -1):

    previous_height, move = parent[i][height]

    path.append(move)

    height = previous_height

path.reverse()

return "".join(path)
\end{lstlisting}
\end{column}
\begin{column}{0.38\textwidth}
\sectiontag{What this chunk does}
\begin{tightitemize}
  \item The only acceptable final height is zero.
  \item If dp[m][0] is unreachable, the answer is IMPOSSIBLE.
  \item Otherwise parent pointers rebuild the moves backward, then reverse them.
\end{tightitemize}
\end{column}
\end{columns}
\end{frame}

\begin{frame}[fragile]{Code: Read Scenarios}
\scriptsize
\sectiontag{Source lines:} \texttt{spiderman.py:142--155}

\vspace{0.08cm}
\begin{columns}[T,totalwidth=\textwidth]
\begin{column}{0.58\textwidth}
\begin{lstlisting}[style=cpcode]
def main():
    input = sys.stdin.buffer.readline

    test_cases = int(input())

    for _ in range(test_cases):
        m = int(input())

        distances = list(map(int, input().split()))

        print(solve(distances))

main()
\end{lstlisting}
\end{column}
\begin{column}{0.38\textwidth}
\sectiontag{What this chunk does}
\begin{tightitemize}
  \item Each test case provides the number of moves and the distance sequence.
  \item The solver returns either a U/D string or IMPOSSIBLE.
\end{tightitemize}
\end{column}
\end{columns}
\end{frame}
% END GENERATED CODE WALKTHROUGH

\begin{frame}{Source}
\small
\sectiontag{Solution file:} \texttt{practice\_contests/day\_3/spiderman.py}

\vspace{0.3cm}
\sectiontag{Problem:} \url{https://open.kattis.com/problems/spiderman}

\vspace{0.45cm}
Use this deck with the source code open beside it. The slides explain the reasoning; the code shows the exact input handling and output formatting.
\end{frame}

\end{document}
